%====================================================
% 01 Introduction
%====================================================

\section{Introduction}

One of the most studied applications of image processing that uses deep learning is handwritten digit recognition. It plays a significant role in many practical applications such as postal mail sorting, automatic cheque reading, automatic form recognition, document digitization and optical character recognition (OCR). The goal of handwritten digit recognition is to correctly classify handwritten numerical digits irrespective of their varying writing habits, size, orientation and stroke width.

In practice, traditional machine learning methods need to extract features from the images prior to classification, which are then used to train a model, and are therefore less effective at handling complex image patterns. Convolutional Neural Networks (CNNs) on the other hand, learn the hierarchy of image features directly from the raw pixel data through multiple layers of convolution and pooling operations. This capability has led CNNs to become one of the most successful deep learning architectures for image classification.

Internationally, the MNIST handwritten digit dataset is today the benchmark for measuring the performance of handwritten digit recognition systems. It is made up of 70,000 gray images of handwritten digits from 0 to 9, each $28 \times 28$ pixels in size. The dataset contains 60,000 training images and 10,000 testing images, making it an excellent source for training and testing CNN model for image classification.

In this assignment, five different CNN architectures were designed and implemented using TensorFlow and Keras in Python. The impact of the different architectural parameters, such as convolution kernel size, number of convolution layers, optimization algorithm and activation function, is analyzed in each model. The models were trained with the MNIST database and tested for their accuracy in the test data, testing loss, confusion matrix analysis, misclassification analysis and custom handwritten digit prediction.

A comparative study was conducted to study the effects of different CNN architectures with respect to classification accuracy, computational complexity and training behaviour. The chosen model that best performed was selected based on its overall recognition accuracy, prediction reliability and computational efficiency.

Lastly, the handwritten digit images, which were made independently, were used for validation of the trained CNN model. The good prediction of these custom images shows that the developed system can generalize from the benchmark ones and recognize handwritten digits in real circumstances.

This assignment gives an all-round knowledge about the classification of images using CNN, evaluation of the model and its comparative analysis and the ability of deep learning techniques to recognize the handwritten digits.